\documentclass[11pt]{article}

\usepackage[margin=1in]{geometry}
\usepackage{amsmath,amssymb,amsthm,mathtools,booktabs,enumitem,microtype,xcolor}
\usepackage[colorlinks=true,linkcolor=blue!50!black,urlcolor=blue!50!black,citecolor=blue!50!black]{hyperref}

\newtheorem{theorem}{Theorem}
\newtheorem{lemma}[theorem]{Lemma}
\newtheorem{proposition}[theorem]{Proposition}
\newtheorem*{corollary}{Corollary}
\theoremstyle{definition}
\newtheorem{definition}{Definition}
\theoremstyle{remark}
\newtheorem*{remark}{Remark}

\DeclareMathOperator{\lca}{lca}
\DeclareMathOperator{\dist}{dist}
\DeclareMathOperator{\htr}{ht}
\DeclareMathOperator{\CF}{CF}
\DeclareMathOperator{\CS}{CS}
\DeclareMathOperator{\labelop}{label}
\DeclareMathOperator{\parent}{parent}
\DeclareMathOperator{\depth}{depth}

\title{Sharp Deterministic Refinement of the OLA--SPR Distance Bound}
\author{}
\date{November 2026}

\begin{document}
\maketitle

\begin{center}
\small
\textbf{Target Reference:} David Harry Richman, Cheng Zhang, and Frederick A. Matsen~IV,
\emph{Vector encoding of phylogenetic trees by ordered leaf attachment},
\href{https://arxiv.org/abs/2503.10169}{arXiv:2503.10169} (v1/v2 Mar 2025, Conjecture~11, \texttt{conj:spr-move-bound});
published in \emph{Bull. Math. Biol.} \textbf{88}:58 (Mar 2026,
\href{https://pmc.ncbi.nlm.nih.gov/articles/PMC12999618/}{PMC12999618}).\\[3pt]
\textbf{Domain:} Mathematical Biology / Phylogenetics and Machine Learning on Tree Spaces.\\[3pt]
\textbf{Formalization:} Combinatorial core formalized in Lean~4 (\texttt{Conjectures/PhyloOLA.lean}).
\end{center}

\medskip

\begin{abstract}
We establish a sharp deterministic bound on the Ordered Leaf Attachment (OLA) Hamming distance across a single rooted subtree prune-and-regraft (rSPR) move. Specifically, for any $T \in \widehat{\mathcal{RBT}}_n$ ($n \ge 2$) and any single-rSPR neighbor $T'$ obtained by pruning at parent $p$ and regrafting below $v_{\mathrm{top}}$,
\[
  d_{\mathrm{OLA}}(T, T') \le \bigl|A_T(p) \cup A_T(v_{\mathrm{top}})\bigr| - 1 \le 2\htr_r(T) - 2 = 2\htr_\rho(T) - 4.
\]
We prove exact worst-case sharpness ($d_{\mathrm{OLA}}(T, T') = 2h - 2$ for every $h \ge 2$ and every $2h \le n \le 2^h$), asymptotic tightness in expectation on caterpillars ($\mathbb{E}[d_{\mathrm{OLA}}(T_n, T')] \ge \frac{n}{24} - O(1) = \Omega(\htr(T_n))$), and the corresponding sharp worst-case NNI bound.
\end{abstract}

\begin{remark}[Novelty and Attribution]
While \href{https://arxiv.org/abs/2503.10169}{arXiv:2503.10169} v1/v2 left the bound as Conjecture~11, the published \emph{Bulletin of Mathematical Biology} version (March 2026) upgraded it to Theorem~6 ($d_{\mathrm{OLA}}(T, T') \le 2\htr_\rho(T)$) using a referee's tripod sketch. Our independent proof yields strictly sharper and exact results not in the published paper:
\begin{enumerate}[leftmargin=2em,itemsep=2pt]
  \item the finer local bound $d_{\mathrm{OLA}}(T, T') \le |A_T(p) \cup A_T(v_{\mathrm{top}})| - 1 \le 2\htr_r(T) - 2 = 2\htr_\rho(T) - 4$ (strictly improving Theorem~6's additive constant by $4$ in the $\htr_\rho$ convention);
  \item exact worst-case sharpness ($d_{\mathrm{OLA}}(T, T') = 2h - 2$ attained for every $h \ge 2$ and every $2h \le n \le 2^h$);
  \item asymptotic tightness in expectation ($\mathbb{E}[d_{\mathrm{OLA}}(T_n, T')] \ge \frac{n}{24} - O(1) = \Omega(\htr(T_n))$ on caterpillars); and
  \item a Lean~4 formalization of the combinatorial core (\texttt{Conjectures/PhyloOLA.lean}).
\end{enumerate}
\end{remark}

\section{Exact Definitions}

Let $\widehat{\mathcal{RBT}}_n$ ($n \ge 2$) denote the set of rooted binary phylogenetic trees on leaf set $L(T) = \{0, 1, \dots, n-1\}$ with a degree-$1$ root $\rho$ and $n - 1$ bifurcating internal vertices $V_{\mathrm{int}}(T)$, whose topmost internal vertex (the unique child of $\rho$) is $r(T)$. Write
\begin{align*}
  \htr_r(T) &:= \max_{\ell \in L(T)} \dist_T\bigl(r(T), \ell\bigr) \ge 1 && \text{(height from the topmost bifurcating node)}, \\
  \htr_\rho(T) &:= \max_{\ell \in L(T)} \dist_T(\rho, \ell) = \htr_r(T) + 1 && \text{(height from $\rho$)}.
\end{align*}
For each $v \in V(T) \setminus \{\rho\}$, let $T_v$ be the rooted leaf-labeled subtree rooted at $v$, and let $L_T(v) \subseteq \{0, \dots, n-1\}$ be its leaf set.

\begin{itemize}[leftmargin=2em,itemsep=4pt]
  \item \emph{Clade-founder and clade-splitter} (Algorithm~1 of arXiv:2503.10169):
  \[
    \CF_T(v) := \min L_T(v),
    \qquad
    \CS_T(v) := \max\bigl(\CF_T(c_1), \CF_T(c_2)\bigr) \quad \bigl(v \in V_{\mathrm{int}}(T)\bigr),
  \]
  where $c_1, c_2$ are the children of $v$. Canonical labels are $\labelop_T(v) := -\CS_T(v)$ for $v \in V_{\mathrm{int}}(T)$ and $\labelop_T(\ell) := \ell$ for leaves $\ell$. By Proposition~4 of arXiv:2503.10169, $\CS_T : V_{\mathrm{int}}(T) \to \{1, \dots, n-1\}$ is a bijection.

  \item \emph{OLA encoding} (Algorithm~2 and Proposition~6): For each $i \in \{1, \dots, n-1\}$, let $T|_i$ be the induced rooted binary tree on leaf set $\{0, 1, \dots, i\}$ (obtained by deleting leaves $> i$ and suppressing degree-$2$ vertices). Then the $i$-th entry $a_i(T)$ of $\Phi(T) = (a_1(T), \dots, a_{n-1}(T))$ is the canonical label in $T|_i$ (equivalently in $T$, by Proposition~4(a)) of the sister of leaf $i$ in $T|_i$.

  \item \emph{OLA distance:} The Hamming distance
  \[
    d_{\mathrm{OLA}}(T, T') := \bigl|\bigl\{i \in \{1, \dots, n-1\} : a_i(T) \neq a_i(T')\bigr\}\bigr|.
  \]

  \item \emph{Rooted SPR (rSPR) move:} Choose $u \in V(T) \setminus \{\rho, r(T)\}$ with parent $p = \parent_T(u) \in V_{\mathrm{int}}(T)$; prune $S = T_u$ by deleting $(p, u)$ and suppressing $p$ to obtain $T \setminus S$; choose an edge $e = (v_{\mathrm{top}}, v_{\mathrm{bot}})$ of $T \setminus S$ (directed away from $\rho$, possibly with $v_{\mathrm{top}} = \rho$); subdivide $e$ by a new internal node $q$ and attach $u$ as its second child to obtain $T'$.
\end{itemize}

\section{Proof of the Main Theorem}

\begin{lemma}[Subtree Locality of OLA Entries]\label{lem:locality}
Let $T, T' \in \widehat{\mathcal{RBT}}_n$. Suppose $w \in V_{\mathrm{int}}(T)$ and $w' \in V_{\mathrm{int}}(T')$ have isomorphic rooted leaf-labeled subtrees $T_w \cong T'_{w'}$. Then $\CS_T(w) = \CS_{T'}(w')$, and setting $i := \CS_T(w)$, we have $a_i(T) = a_i(T')$.
\end{lemma}

\begin{proof}
Let $\varphi : V(T_w) \to V(T'_{w'})$ be the rooted leaf-label-preserving isomorphism with $\varphi(w) = w'$. For every $x \in V(T_w)$, $L_{T'}(\varphi(x)) = L_T(x)$, so $\CF_{T'}(\varphi(x)) = \CF_T(x)$ and $\labelop_{T'}(\varphi(x)) = \labelop_T(x)$. In particular, $i := \CS_T(w) = \CS_{T'}(w') \in \{1, \dots, n-1\}$.

Order the children $c_1, c_2$ of $w$ in $T$ so that $\CF_T(c_1) < \CF_T(c_2) = i$, and set $A := \{0, 1, \dots, i\}$. Then $L_T(c_2) \cap A = \{i\}$ and $L_T(c_1) \cap A \neq \emptyset$. A vertex $x \in V(T) \setminus \{\rho\}$ survives in $T|_i$ if and only if $x \in A$ or $x$ is an internal node both of whose children have clades intersecting $A$. Thus $w$ survives in $T|_i$: its child on the $c_2$-side in $T|_i$ is leaf $i$, and its child on the $c_1$-side in $T|_i$ is the highest surviving vertex in $T_{c_1}$, namely
\[
  z := \lca_{T_{c_1}}\!\bigl(L_T(c_1) \cap A\bigr).
\]
By canonical-label stability under restriction (Proposition~4(a) and Proposition~6 of arXiv:2503.10169), $a_i(T) = \labelop_{T|_i}(z) = \labelop_T(z)$. Applying the same argument to $w'$ in $T'$ gives the sister of leaf $i$ in $T'|_i$ as $z' = \lca_{T'_{\varphi(c_1)}}(L_{T'}(\varphi(c_1)) \cap A) = \varphi(z)$, hence
\[
  a_i(T') = \labelop_{T'}(\varphi(z)) = \labelop_T(z) = a_i(T). \qedhere
\]
\end{proof}

\begin{lemma}[Upward-Closed Change Set Bound]\label{lem:upward}
Let $T, T' \in \widehat{\mathcal{RBT}}_n$. Suppose $B \subseteq V_{\mathrm{int}}(T)$ is a nonempty upward-closed subset of internal nodes of $T$ (that is, if $w \in B$ and $v \in V_{\mathrm{int}}(T)$ is an ancestor of $w$, then $v \in B$) such that for every $w \in V_{\mathrm{int}}(T) \setminus B$, there exists $w' \in V_{\mathrm{int}}(T')$ with $T'_{w'} \cong T_w$. Then
\[
  d_{\mathrm{OLA}}(T, T') \le |B| - 1.
\]
\end{lemma}

\begin{proof}
Let $J := \{\CS_T(w) : w \in B\} \subseteq \{1, \dots, n-1\}$. By Proposition~4 of arXiv:2503.10169, $|J| = |B| \ge 1$. For every $i \in \{1, \dots, n-1\} \setminus J$, the unique $w \in V_{\mathrm{int}}(T)$ with $\CS_T(w) = i$ lies outside $B$, so Lemma~\ref{lem:locality} gives $a_i(T) = a_i(T')$.

Let $j^* := \min J$. We show that $a_{j^*}(T) = a_{j^*}(T')$ as well:
\begin{itemize}[leftmargin=2em,itemsep=3pt]
  \item If $j^* = 1$, then $a_1(T) = 0 = a_1(T')$.
  \item If $j^* \ge 2$, let $v^* := \lca_T(\{0, 1, \dots, j^* - 1\}) \in V_{\mathrm{int}}(T)$. Both children of $v^*$ intersect $\{0, \dots, j^* - 1\}$, so $\CS_T(v^*) \le j^* - 1 < j^* = \min J$, forcing $v^* \notin B$. Hence there exists $v^{*\prime} \in V_{\mathrm{int}}(T')$ with $T'_{v^{*\prime}} \cong T_{v^*}$.

  We claim $j^* \notin L_T(v^*)$. Otherwise, let $w^* \in B$ be the unique internal node with $\CS_T(w^*) = j^*$; then $j^* \in L_T(w^*)$, so $v^*$ and $w^*$ are comparable in the ancestor order:
  \begin{enumerate}[label=\textup{(\roman*)},leftmargin=2em,itemsep=2pt]
    \item $v^*$ cannot be an ancestor of $w^*$, since $w^* \in B$ and upward-closure of $B$ would imply $v^* \in B$, contradicting $v^* \notin B$.
    \item $w^*$ cannot be a proper ancestor of $v^*$, since then the child $c$ of $w^*$ containing $v^*$ would satisfy $L_T(c) \supseteq L_T(v^*) \supseteq \{0, \dots, j^*\}$, forcing $\CF_T(c) = 0$ and $\CF_T(c') \ge j^* + 1$ for the other child $c'$, contradicting $\CS_T(w^*) = j^*$.
  \end{enumerate}
  Hence $j^* \notin L_T(v^*) = L_{T'}(v^{*\prime})$. Therefore $L_T(v^*) \cap \{0, \dots, j^*\} = \{0, \dots, j^* - 1\}$, so the two children of the topmost node of $T|_{j^*}$ are $v^*$ and leaf $j^*$, giving $a_{j^*}(T) = \labelop_T(v^*) = -\CS_T(v^*)$. Since $T'_{v^{*\prime}} \cong T_{v^*}$, the same holds in $T'|_{j^*}$, yielding $a_{j^*}(T') = -\CS_{T'}(v^{*\prime}) = -\CS_T(v^*) = a_{j^*}(T)$.
\end{itemize}
Thus $a_i(T) = a_i(T')$ for all $i \notin J \setminus \{j^*\}$, so $d_{\mathrm{OLA}}(T, T') \le |J| - 1 = |B| - 1$.
\end{proof}

\begin{theorem}[Sharp Deterministic rSPR Bound; Proof of Conjecture~11]\label{thm:rspr}
Let $T \in \widehat{\mathcal{RBT}}_n$ ($n \ge 2$), and let $T'$ be obtained from $T$ by a single rSPR move pruning $S = T_u$ ($p = \parent_T(u) \in V_{\mathrm{int}}(T)$) and regrafting onto edge $e = (v_{\mathrm{top}}, v_{\mathrm{bot}})$ of $T \setminus S$. For $x \in V_{\mathrm{int}}(T)$, let $A_T(x)$ be the set of internal ancestors of $x$ in $T$ (including $x$), and set $A_T(\rho) := \emptyset$. Then
\[
  d_{\mathrm{OLA}}(T, T') \le \bigl|A_T(p) \cup A_T(v_{\mathrm{top}})\bigr| - 1 \le 2\htr_r(T) - 2 = 2\htr_\rho(T) - 4.
\]
In particular, by symmetry of rSPR, $d_{\mathrm{OLA}}(T, T') \le 2\min\bigl(\htr_r(T), \htr_r(T')\bigr) - 2$, and for any probability distribution on single-SPR neighbors $T'$ of $T$,
\[
  \mathbb{E}\bigl[d_{\mathrm{OLA}}(T, T')\bigr] \le 2\htr_r(T) - 2 = O\bigl(\htr(T)\bigr).
\]
\end{theorem}

\begin{proof}
Note that $v_{\mathrm{top}} \in \{\rho\} \cup (V_{\mathrm{int}}(T) \setminus \{p\})$. Set $B := A_T(p) \cup A_T(v_{\mathrm{top}}) \subseteq V_{\mathrm{int}}(T)$, which is nonempty ($p \in B$) and upward-closed in $T$. For every $w \in V_{\mathrm{int}}(T) \setminus B$:
\begin{itemize}[leftmargin=2em,itemsep=2pt]
  \item $w \neq p$, so $w$ is not suppressed when forming $T \setminus S$.
  \item If $w \in V_{\mathrm{int}}(S)$, then $T'_u \cong T_u$, hence $T'_w \cong T_w$.
  \item If $w \notin V_{\mathrm{int}}(S)$, then $w \notin A_T(p)$ implies $T_w$ is disjoint from $\{p\} \cup V(S)$, so $(T \setminus S)_w \cong T_w$. Since the internal ancestors of the new node $q$ in $T'$ (other than $q$) are $A_T(v_{\mathrm{top}}) \setminus \{p\} \subseteq B$, $w$ is not an ancestor of $q$ in $T'$, so subdividing $e$ leaves the subtree below $w$ untouched: $T'_w \cong (T \setminus S)_w \cong T_w$.
\end{itemize}
By Lemma~\ref{lem:upward}, $d_{\mathrm{OLA}}(T, T') \le |B| - 1$. Every $x \in V_{\mathrm{int}}(T)$ has a leaf strictly below it, so $|A_T(x)| = \depth_r(x) + 1 \le \htr_r(T)$.
\begin{itemize}[leftmargin=2em,itemsep=2pt]
  \item If $v_{\mathrm{top}} = \rho$, then $|B| - 1 = |A_T(p)| - 1 \le \htr_r(T) - 1 \le 2\htr_r(T) - 2$.
  \item If $v_{\mathrm{top}} \in V_{\mathrm{int}}(T) \setminus \{p\}$, then $r(T) \in A_T(p) \cap A_T(v_{\mathrm{top}})$, so $|B| \le |A_T(p)| + |A_T(v_{\mathrm{top}})| - 1 \le 2\htr_r(T) - 1$, hence $|B| - 1 \le 2\htr_r(T) - 2$.
\end{itemize}
\end{proof}

\section{Exact Sharpness and Corollaries}

\begin{proposition}[Worst-Case Exact Sharpness]\label{prop:sharp}
For every $h \ge 2$ and every $n$ satisfying $2h \le n \le 2^h$, there exists $T \in \widehat{\mathcal{RBT}}_n$ with $\htr_r(T) = h$ and a single rSPR neighbor $T'$ with $d_{\mathrm{OLA}}(T, T') = 2h - 2$.
\end{proposition}

\begin{proof}
Any rooted binary tree of height $h$ has at most $2^h$ leaves, so $n \le 2^h$ is necessary. For $n = 2h$, let $T$ have topmost node $r$ with two height-$(h-1)$ caterpillar children: left spine $L_1, \dots, L_{h-1}$ attaching leaves $2, 3, \dots, h$ with bottom leaf $0$, and right spine $R_1, \dots, R_{h-1}$ attaching leaves $h+1, \dots, 2h-1$ with bottom leaf $1$. Then $\Phi(T) = (0^h, 1^{h-1})$. Pruning leaf $0$ and regrafting it onto the parent edge of leaf $1$ gives $T'$ with
\[
  \Phi(T') = \bigl(0,\; -1,\; 2,\; 3,\; \dots,\; h-1,\; (-1)^{h-1}\bigr),
\]
which differs from $\Phi(T)$ in all $2h - 2$ coordinates $i = 2, \dots, 2h-1$. For $2h < n \le 2^h$, expanding side leaves $\{2, \dots, 2h-1\}$ into binary subtrees on the extra labels $\{2h, \dots, n-1\}$ while keeping height $\le h$ (and leaving leaves $0, 1$ and the spine structure intact) preserves all $2h - 2$ changed entries at coordinates $2, \dots, 2h-1$, while every new internal node lies outside $B$ and has unchanged entry by Lemma~\ref{lem:locality}, giving $d_{\mathrm{OLA}}(T, T') = 2h - 2$.
\end{proof}

\begin{proposition}[Asymptotic Tightness in Expectation]\label{prop:expectation}
For the caterpillar $T_n = \Psi(0, \dots, 0) \in \widehat{\mathcal{RBT}}_n$ ($\htr_r(T_n) = n - 1$) and $T'$ drawn uniformly from $\mathcal{N}_{\mathrm{SPR}}(T_n)$,
\[
  \mathbb{E}\bigl[d_{\mathrm{OLA}}(T_n, T')\bigr] \ge \frac{n}{24} - O(1) = \Omega\bigl(\htr(T_n)\bigr).
\]
\end{proposition}

\begin{proof}
For $1 \le j \le k - 2 \le n - 4$, pruning leaf $j$ and regrafting onto the parent edge of internal node $-(k+1)$ produces pairwise distinct caterpillars $T'(j, k)$ (uniquely identified by their single descent $(k, j)$ in the top-to-bottom leaf peel sequence) with $a_m(T'(j, k)) = -j \neq 0 = a_m(T_n)$ for all $m \in \{j+1, \dots, k\}$, so $d_{\mathrm{OLA}}(T_n, T'(j, k)) \ge k - j$. Since $|\mathcal{N}_{\mathrm{SPR}}(T_n)| < 4n^2$, summing over $d = k - j \in \{2, \dots, n-3\}$ yields
\[
  \mathbb{E}\bigl[d_{\mathrm{OLA}}(T_n, T')\bigr] \ge \frac{1}{4n^2} \sum_{d=2}^{n-3} d(n - 2 - d) = \frac{n}{24} - O(1). \qedhere
\]
\end{proof}

\begin{corollary}[Sharp Worst-Case NNI Bound]
For a nearest-neighbor interchange (NNI) move on internal edge $(v_1, v_2)$, $A_T(p) \cup A_T(v_{\mathrm{top}}) = A_T(v_2)$, so Theorem~\ref{thm:rspr} yields
\[
  d_{\mathrm{OLA}}(T, T') \le \depth_r(v_2) \le \htr_r(T) - 1.
\]
\end{corollary}

\end{document}
